\chapter{Kiến thức và công nghệ nền tảng}
\label{chap:coso}

Chương này trình bày kiến thức Farmery vận dụng: đặc điểm mặt hàng nông sản, bối cảnh chính sách và pháp lý, các lý thuyết về mô hình thương mại điện tử, truy xuất nguồn gốc, uy tín và các bài toán AI. Các nền tảng và nghiên cứu cụ thể được khảo sát riêng tại Chương~\ref{chap:lienquan}.

\section{Đặc điểm hàng nông sản}
\label{sec:dacdiemnongsan}

\subsection{Phạm vi khái niệm nông sản trong đồ án}
\label{subsec:dinhnghianongsan}

Pháp luật Việt Nam không có định nghĩa duy nhất về ``nông sản''. Đồ án dựa vào ba căn cứ: Luật An toàn thực phẩm 55/2010 phân biệt \emph{thực phẩm tươi sống}, \emph{sơ chế} và \emph{chế biến} \parencite{luat_attp_2010}; Nghị định 15/2018 buộc \emph{thực phẩm chế biến bao gói sẵn} phải tự công bố sản phẩm, còn nông sản tươi hoặc sơ chế thì không \parencite{nd15_2018}; Luật Trồng trọt 31/2018 nhìn chuỗi giá trị là liên tục từ thu hoạch, sơ chế, bảo quản đến thương mại \parencite{luat_trongtrot_2018}. Từ đó hàng hóa chia bốn nhóm:
\begin{itemize}
    \item \textbf{Nhóm 1 --- Nông sản tươi} (rau, củ, quả, gạo, hạt) và \textbf{Nhóm 2 --- sơ chế, bảo quản đơn giản} (sấy, cấp đông, hút chân không, không phụ gia): \emph{trong phạm vi}, là trục chính.
    \item \textbf{Nhóm 3 --- Chế biến sâu} (nước ép, mứt, trà, tinh dầu; phần lớn sản phẩm OCOP): \emph{trong phạm vi có điều kiện}, nhà cung cấp phải có bản tự công bố sản phẩm và giấy chứng nhận cơ sở đủ điều kiện an toàn thực phẩm.
    \item \textbf{Nhóm 4 --- Ngoài phạm vi:} thủy sản sống, thịt tươi, động vật sống, thực phẩm chức năng, vật tư nông nghiệp và hàng phi thực phẩm.
\end{itemize}
Ranh giới này là căn cứ xây dựng danh mục sản phẩm đóng và danh sách loại trừ ở khâu đăng bán (mục~\ref{sec:quytrinhnghiepvu}).

\subsection{Nông sản khác hàng hóa thông thường ở điểm nào}
\label{subsec:sosanhnongsan}

So với hàng công nghiệp và hàng đóng gói mà các sàn phổ thông phục vụ, nông sản có bốn khác biệt chi phối thiết kế:
\begin{itemize}
    \item \textbf{Mau hỏng:} hạn dùng tính bằng ngày, tổn thất sau thu hoạch 20--40\% (Chương~\ref{chap:gioithieu}); tồn kho phải quản lý theo lô kèm hạn dùng, xuất lô hết hạn trước, thời gian giao là ràng buộc chứ không phải tiện ích.
    \item \textbf{Nguồn cung theo mùa vụ:} giá và sản lượng biến động mạnh theo thời điểm, cần dự báo thay vì chỉ ghi nhận.
    \item \textbf{Chất lượng không đồng nhất giữa các lô:} dữ liệu mô tả phải gắn theo từng lô, và hệ thống phải hỗ trợ nhiều đơn vị tính.
    \item \textbf{Người mua không tự kiểm chứng được chất lượng,} kể cả sau khi dùng: đúng loại hàng mà lý thuyết bất cân xứng thông tin mô tả (mục~\ref{subsec:reputation}); truy xuất nguồn gốc và chứng nhận bên thứ ba vì thế là bắt buộc, và kiểm duyệt không thể chỉ dựa vào tự động hóa.
\end{itemize}

\subsection{Chứng nhận, giấy tờ và truy xuất nguồn gốc quốc gia}
\label{subsec:chungnhan}

Một lô nông sản ra thị trường cần hai loại hồ sơ. \textbf{Chứng nhận cách thức sản xuất,} cấp cho cơ sở: nhóm thực hành nông nghiệp tốt (GlobalG.A.P. \parencite{globalgap_ifa}, VietGAP), nhóm hữu cơ (USDA Organic \parencite{usda_nop}, EU Organic \parencite{eu_organic_2018848}, JAS) và nhóm an toàn thực phẩm cho chế biến (HACCP, ISO 22000, FSSC 22000, BRCGS, IFS --- bốn chuẩn sau nằm trong danh sách GFSI công nhận \parencite{gfsi_cpo}). \textbf{Giấy tờ theo từng lô,} chủ yếu cho xuất khẩu: kiểm dịch thực vật theo ISPM 12 \parencite{ippc_ispm12}, chứng nhận xuất xứ, kết quả phân tích dư lượng, kèm mã số vùng trồng và mã số cơ sở đóng gói theo Nghị định 38/2026/NĐ-CP \parencite{ppd_masovungtrong}. Trong đồ án, hệ thống hỗ trợ tải lên và kiểm duyệt nhóm chứng nhận cách thức sản xuất; giấy tờ theo lô cho xuất khẩu nằm ngoài phạm vi (mục~\ref{sec:phamvi}).

Từ 1/7/2026, Bộ Nông nghiệp và Môi trường vận hành chính thức hệ thống truy xuất nguồn gốc nông, lâm, thủy sản của Nhà nước \parencite{traceviet_2026}; cách tổ chức dữ liệu (nông hộ, vùng trồng, lô hàng, cơ sở đóng gói) trùng với nhật ký canh tác theo lô của Farmery nên được dùng làm tham chiếu thiết kế lược đồ (mục~\ref{subsec:dacta-csdl}).

\label{subsec:nhatkycanhtac-xuatkhau}
\noindent\textbf{Vì sao xuất khẩu nằm ngoài phạm vi.} Lệnh 248 của Hải quan Trung Quốc buộc cơ sở thực phẩm nước ngoài đăng ký mã số, và các nhóm rủi ro cao như rau quả tươi phải do cơ quan thẩm quyền nước xuất khẩu tiến cử \parencite{gacc_decree248}; Điều 18 Quy định (EC) 178/2002 buộc mỗi tác nhân xác định nhà cung cấp liền trước và khách hàng liền sau \parencite{eu_gfl_178_2002}, với thời hạn lưu hồ sơ nằm ở hướng dẫn của Ủy ban châu Âu \parencite{ec_gfl_guidance}. Đối chiếu cho thấy khoảng cách không phải ``thêm vài trường dữ liệu'': xác minh chứng nhận phải đi qua đầu mối Nhà nước, định danh vùng trồng phải là mã số chính thức, và chuỗi truy xuất giãn từ hai lên năm sáu bên độc lập. Khối lượng đó vượt một đồ án một học kỳ; điều kiện mở rộng nêu tại Chương~\ref{chap:ketluan}.

\section{Bối cảnh chính sách và tiêu chuẩn chất lượng nông sản}
\label{sec:boicanhchinhsach}

Nông nghiệp là lĩnh vực ưu tiên chuyển đổi số theo Quyết định 749/QĐ-TTg \parencite{qd749_2020} và Nghị quyết 57-NQ/TW \parencite{nq57_tw}; Quyết định 5943/QĐ-BNNMT (31/12/2025) cụ thể hóa bằng cơ sở dữ liệu dùng chung, nguyên tắc khai báo một lần và ứng dụng dữ liệu lớn, AI trong giám sát, dự báo sản xuất \parencite{botnmt_2025, qd5943_2025}. Chương trình OCOP giai đoạn 2026--2030 chuyển trọng tâm ``từ số lượng sang chất lượng'', đẩy mạnh truy xuất nguồn gốc và tiêu thụ qua thương mại điện tử \parencite{danviet_2024_ocop, vietnamnet_2025_ocop} --- vừa thuận lợi cho định hướng của Farmery, vừa là nguồn cung sản phẩm có thương hiệu và câu chuyện xuất xứ.

Về tiêu chuẩn, \textbf{VietGAP} (khoảng 70 tiêu chí) phù hợp chuỗi nội địa; \textbf{GlobalGAP} (252 tiêu chuẩn, mã định danh GGN 13 chữ số) là điều kiện vào thị trường khó tính. Cơ sở đạt GlobalGAP đã vượt chuẩn VietGAP, còn VietGAP thường là bước đệm \parencite{vietgap_globalgap}. Farmery vì vậy kiểm duyệt chứng nhận theo hai cấp (nội địa/VietGAP, xuất khẩu/GlobalGAP) tương ứng khách lẻ và khách sỉ.

\textbf{Khung pháp lý.} Theo Luật Thương mại điện tử 2025, Farmery đồng thời là \emph{nền tảng kinh doanh trực tiếp} (bán lẻ) và \emph{nền tảng trung gian} (bán sỉ cho nhà cung cấp), nên chịu cả hai nhóm nghĩa vụ \parencite{luattmdt2025}. Bảng~\ref{tab:phaply} tóm tắt các văn bản chính và tác động tới thiết kế.

\begingroup
\tblsetup\tblzebra
\begin{longtable}{L{3.50cm} L{6.70cm} L{4.70cm}}
\caption{Văn bản pháp luật, nghĩa vụ và tác động tới thiết kế Farmery}
\label{tab:phaply} \\
\toprule
\rowcolor{tblheadbg}
\textbf{Văn bản} & \textbf{Nghĩa vụ chính đối với Farmery} & \textbf{Tác động thiết kế} \\
\midrule
\endfirsthead
\toprule
\rowcolor{tblheadbg}
\textbf{Văn bản} & \textbf{Nghĩa vụ chính đối với Farmery} & \textbf{Tác động thiết kế} \\
\midrule
\endhead
\bottomrule
\endfoot
Luật Thương mại điện tử 122/2025 và Nghị định 248/2026, hiệu lực 1/7/2026 \parencite{luattmdt2025, nd248_2026} & Công khai chính sách và để người dùng đồng ý trước khi mở tài khoản; hiển thị đủ tiền hàng, thuế, phí giao trước khi đặt; nền tảng trung gian xác thực người bán, kiểm duyệt nội dung, báo trước 5 ngày khi hạn chế tài khoản; lưu dữ liệu hợp đồng ít nhất 3 năm; bộ chính sách có nội dung tối thiểu theo từng loại & FR1.4, FR1.6, FR5.2, FR10.1; NFR8.1, NFR8.5; Phụ lục B \\
Luật Bảo vệ quyền lợi người tiêu dùng 19/2023 \parencite{luatbvqlntd2023} & Giao dịch từ xa: cho xem, xác nhận và tải lại hợp đồng; người tiêu dùng chỉ được đơn phương chấm dứt khi người bán cung cấp thông tin sai hoặc thiếu --- luật không đặt quyền ``đổi ý'' chung & FR5.4, FR7.2; chính sách không đổi ý với hàng tươi hợp lệ khi công khai trước \\
Luật Bảo vệ dữ liệu cá nhân 91/2025 và Nghị định 356/2025 \parencite{luatbvdlcn2025, nd356_2025} & Xử lý dữ liệu có đồng ý, đúng mục đích; quyền xem, sửa, xóa của chủ thể dữ liệu & FR1.6, NFR8.2, NFR3.2 \\
Thuế và hóa đơn: Luật GTGT 48/2024 và 149/2025, Nghị định 117/2025, 68/2026, 70/2025 \parencite{luatgtgt2024, luat149_2025, nd117_2025, nd68_2026, nd70_2025} & Nông dân tự bán: không chịu thuế; nông sản chưa chế biến ở khâu thương mại: 5\%; sàn có chức năng đặt hàng và thanh toán khấu trừ, nộp thuế thay người bán là hộ, cá nhân; lập hóa đơn điện tử đúng nội dung và thời điểm & FR2.4, FR7.2, FR7.4 (chỉ thiết kế, áp cho sàn sỉ), NFR8.3, NFR8.4 \\
Nghị định 37/2026 về nhãn \parencite{nd37_2026} & Thực phẩm tươi không bao bì bán trực tiếp không phải ghi nhãn; đơn vị đã đóng gói phải có nhãn đủ nội dung; nền tảng lưu dữ liệu nhãn 5 năm; sản phẩm bị thu hồi phải hiện cảnh báo & FR13.4, FR13.6, NFR8.5 \\
Luật An toàn thực phẩm 55/2010 và Nghị định 15/2018 \parencite{luat_attp_2010, nd15_2018} & Cơ sở sơ chế, đóng gói phải có giấy chứng nhận cơ sở đủ điều kiện, trừ cơ sở nhỏ lẻ hoặc đã có HACCP, ISO 22000; nông hộ sản xuất ban đầu nhỏ lẻ không thuộc diện cấp & Điều kiện vận hành kho sơ chế, không phải yêu cầu phần mềm \\
Luật Giao dịch điện tử 20/2023 \parencite{luatgddt2023} & Chữ ký số tương đương chữ ký tay; hình thức xác nhận khác theo pháp luật liên quan & FR6.3 gọi là ``xác nhận giao kết''; ký số ở hướng phát triển \\
\end{longtable}
\endgroup

Ba điểm được viết dè dặt: (i) hệ thống chỉ lưu loại người bán và để nhân sự xác nhận việc khấu trừ thuế với nông dân tự sản xuất, không tự suy ra tỷ lệ; (ii) Nghị định 46/2026/NĐ-CP thay Nghị định 15/2018 đang bị tạm ngưng hiệu lực \parencite{nq15_2026}, điều kiện kho sơ chế có thể đổi theo văn bản mới; (iii) ngưỡng doanh thu không chịu thuế \parencite{nd141_2026} để ở cấu hình chứ không ghi cứng vào thiết kế.

\section{Cơ sở lý thuyết nền tảng}
\label{sec:lythuyet}

\subsection{Mô hình thương mại điện tử và chuỗi cung ứng ngắn}
\label{subsec:haitruc}

Thương mại điện tử được phân loại theo hai tiêu chí độc lập. Tiêu chí thứ nhất là \textbf{bản chất các bên:} B2B (đơn hàng lớn theo hợp đồng), B2C (doanh nghiệp bán cho người tiêu dùng), B2B2C (ba mắt xích). Tiêu chí thứ hai, quyết định với Farmery, là \textbf{ai sở hữu hàng trong lúc giao dịch:} mô hình \emph{third-party} (3P) để người bán giữ hàng tới khi giao, nền tảng chỉ trung gian thu hoa hồng; mô hình \emph{first-party} (1P) để nền tảng mua vào, tự định giá, bán lại và gánh rủi ro tồn kho. Cùng là B2B2C, Shopee theo 3P còn mảng bán lẻ của Amazon theo 1P. Trục thứ nhất quyết định nhóm người dùng cần phục vụ; trục thứ hai quyết định ai chịu rủi ro tài chính và trách nhiệm chất lượng. Farmery vận hành bán sỉ theo B2B-3P và bán lẻ theo B2B2C-1P (lý do ở mục~\ref{sec:businessmodelcanvas}).

Chuỗi cung ứng thực phẩm ngắn (SFSC) là chuỗi có ít khâu trung gian và quan hệ gần gũi giữa nhà sản xuất và người tiêu dùng \parencite{tapchicongsan2018}: tăng minh bạch, niềm tin và lợi nhuận cho nhà sản xuất, nhưng khó mở rộng quy mô và chi phí đơn vị cao. Liên quan là \emph{disintermediation} --- công nghệ số cho nhà sản xuất tiếp cận trực tiếp người mua cuối, như Dell bán qua web \parencite{britannica_disintermediation}. Farmery là một nỗ lực disintermediation có kiểm soát: bỏ các tầng thương mại thuần túy hưởng chênh lệch giá, nhưng vẫn hợp tác đơn vị vận chuyển và giữ lợi ích quy mô của chuỗi dài nhờ nền tảng số.

\subsection{Truy xuất nguồn gốc: chuẩn GS1 và vai trò của blockchain}
\label{subsec:gs1blockchain}

GS1 xây dựng chuẩn định danh dùng chung cho chuỗi cung ứng: GTIN cho sản phẩm, GLN cho bên và địa điểm, SSCC cho đơn vị hậu cần, cho phép truy xuất thông tin lô khi cần thu hồi hoặc xác minh xuất xứ \parencite{gs1_keys}. Farmery áp dụng nguyên lý này: mỗi lô có mã QR định danh gắn với dữ liệu của lô.

Blockchain bảo đảm dữ liệu không bị sửa \emph{sau khi ghi} nhưng không xác minh dữ liệu đầu vào đúng hay sai; nông dân và cơ sở nhỏ vẫn nhập thủ công nên gian lận tại nguồn không bị phát hiện \parencite{bath_blockchain}. Trường hợp thành công điển hình --- Walmart với IBM Food Trust rút thời gian truy vết lô xoài từ bảy ngày xuống 2,2 giây \parencite{walmart_ibm_foodtrust} --- có tiền đề là Walmart đủ lớn để buộc nhà cung cấp nhập liệu đúng chuẩn, điều một sàn mới với nhiều nhà cung cấp nhỏ lẻ không có. Chi phí hạ tầng không phải căn cứ: một nghiên cứu ở Thái Lan cho thấy blockchain rẻ hơn khoảng 43\% so với cách làm tập trung \parencite{thai_broiler_blockchain_2025}. Căn cứ thật là rào cản áp dụng: bảo mật và độ trưởng thành công nghệ, quyền riêng tư, thiếu hạ tầng số \parencite{perishable_blockchain_barriers_2025}; rào cản thứ ba trùng với hạn chế kỹ năng số của nhóm nhà cung cấp (mục~\ref{sec:doituongnguoidung}). Farmery vì vậy lưu nhật ký canh tác trên cơ sở dữ liệu tập trung theo nguyên lý GS1, kèm cơ chế bất biến sau khi công khai (NFR9.1).

Hệ thống truy xuất quốc gia (TraceViet) thí điểm từ 26/12/2025 và vận hành chính thức từ 1/7/2026 \parencite{tuoitre_traceviet_2025, mae_traceviet_2026}. Chuẩn kết nối do Cục Chuyển đổi số công bố \parencite{traceviet_chuanketnoi} cho thấy: nền tảng tham gia với tư cách \emph{đơn vị cung cấp giải pháp}, phải đăng ký và qua kiểm thử kỹ thuật; chuẩn chỉ mô tả chiều \emph{đẩy} dữ liệu bằng tệp XML ký số, không có giao diện tra cứu cho bên ngoài; mã truy xuất chính thức do hệ thống của Bộ cấp sau phê duyệt. Farmery thiết kế lớp tích hợp theo chuẩn đã công bố (chỉ thiết kế); mỗi lô vẫn có mã QR riêng sinh ngay khi được duyệt nội bộ và việc mở bán không phụ thuộc đồng bộ lên hệ thống của Bộ.

\subsection{Hệ thống đánh giá, uy tín và niềm tin}
\label{subsec:reputation}

Mô hình ``The Market for Lemons'' của Akerlof \parencite{akerlof1970lemons} cho thấy khi người bán biết rõ chất lượng hơn người mua và người mua không kiểm chứng được, hàng kém đào thải hàng tốt (lựa chọn đối nghịch): người bán chất lượng cao không chứng minh được giá trị để nhận mức giá tương xứng. Với sàn nông sản, đó là bài toán cốt lõi vì người tiêu dùng khó phân biệt nông sản ``sạch'' thật với sản phẩm chỉ gắn mác marketing. Hệ thống uy tín giảm bất cân xứng bằng cách tích lũy phản hồi (điểm, bình luận, tỷ lệ hoàn thành đơn) thành chỉ số công khai cho từng người bán \parencite{springer_article}, dù có hạn chế như ngại đánh giá tiêu cực hay đánh giá giả. Với Farmery, đánh giá chủ quan được bổ sung bằng hai lớp khách quan: dữ liệu truy xuất và chứng nhận đã xác minh.

\subsection{Công nghệ nền tảng, mô hình doanh thu và vận hành kho}
\label{subsec:businessmodel}

\textbf{Kiến trúc và phân quyền.} Sàn đa vai trò gồm lớp giao diện theo từng nhóm người dùng, lớp API và nghiệp vụ, lớp dữ liệu (cơ sở dữ liệu quan hệ kèm object storage cho ảnh, chứng nhận), tích hợp bên thứ ba qua API. Kiểm soát truy cập theo \emph{RBAC} \parencite{sandhu1996rbac}: quyền cấp qua vai trò thay vì cho từng người dùng, thực thi nguyên tắc đặc quyền tối thiểu.

\textbf{Thanh toán.} Trung gian thanh toán là tổ chức được cấp phép thu hộ, chi hộ (VNPay, Momo). Với giao dịch B2B giá trị lớn có thể dùng \emph{escrow}: tiền người mua được bên trung gian giữ và chỉ giải ngân khi điều kiện đã thỏa thuận được đáp ứng.

\textbf{Mô hình doanh thu.} Theo Business Model Canvas \parencite{osterwalder2010bmc}, doanh thu marketplace thường gồm hoa hồng giao dịch, phí niêm yết, dịch vụ giá trị gia tăng và quảng cáo. Giai đoạn đầu ưu tiên hoa hồng vì gắn với giá trị giao dịch thực và ít cản người bán mới; mô hình của Farmery phân tích tại mục~\ref{sec:businessmodelcanvas}.

\label{subsec:khobai}
\textbf{Tồn kho và vận chuyển.} Hàng dễ hư áp dụng FEFO (hết hạn trước xuất trước) và quản lý theo lô kèm cảnh báo sớm khi lô sắp hết hạn, nhằm giảm hao hụt 20--40\% đã nêu ở Chương~\ref{chap:gioithieu}. Nhiều sàn thuê ngoài vận chuyển (3PL) để tập trung vào công nghệ, đổi lại phụ thuộc chất lượng đối tác. Farmery thuê ngoài phương tiện nhưng tự giữ các khâu 3PL phục vụ kém nhất với hàng tươi: nghiệm thu, đóng gói, dán tem tại kho sơ chế và điều phối xe lạnh chặng dài; chặng cuối nội thành giao cho 3PL. Đội giao hàng riêng chỉ đáng khi mật độ đơn đủ lớn nên để ở hướng phát triển (ranh giới tại mục~\ref{sec:phamvi}).

\section{Ứng dụng AI trong thương mại điện tử nông sản}

Farmery xét bốn nhóm bài toán AI hỗ trợ ra quyết định, trong đó ba nhóm là mô-đun hiện thực (tìm kiếm ngữ nghĩa, gợi ý sản phẩm, dự báo giá) và giám sát bất thường dùng luật ngưỡng; thiết kế chi tiết ở mục~\ref{sec:thietkeAI}, hiện thực ở mục~\ref{sec:hienthucAI}.

\textbf{Tìm kiếm theo ngữ nghĩa.}\label{subsec:semanticsearch} Tìm theo từ khóa lệch với cách người dùng gọi tên nông sản (``rau ăn lẩu'' so với ``cải cúc''). Tìm kiếm ngữ nghĩa biểu diễn truy vấn và mô tả thành vector trong cùng không gian, xếp theo độ gần \parencite{reimers2019sbert}; vì yếu ở khớp chính xác tên riêng, mã, vùng trồng nên dùng \emph{tìm kiếm lai} hợp nhất hai danh sách bằng reciprocal rank fusion \parencite{cormack2009rrf}. Với tiếng Việt cần mô hình nhúng đa ngôn ngữ như BGE-M3 \parencite{chen2024bgem3}; chất lượng đo bằng nDCG \parencite{jarvelin2002ndcg}.

\textbf{Gợi ý sản phẩm.} Hai hướng phổ biến là collaborative filtering (theo tương đồng hành vi người dùng) và content-based filtering (theo đặc điểm sản phẩm), thường kết hợp để giảm vấn đề ``cold start'' \parencite{su2009collaborative}. Farmery tận dụng đặc thù dữ liệu nông sản (mùa vụ, vùng trồng, chứng nhận) để gợi ý sản phẩm phù hợp thời điểm.

\textbf{Dự báo giá.} Đây là bài toán chuỗi thời gian, kết hợp mô hình thống kê (ARIMA, SARIMA) với học máy khi dữ liệu đủ lớn \parencite{hyndman2021forecasting}; giúp nhân viên thu mua quyết định mua lô nào, giá nào và đặt giá bán lẻ còn lãi sau hao hụt. Hướng này đã được kiểm chứng: Ninjacart dự báo nhu cầu bằng học máy \parencite{ninjacart}, Farm2Market công bố độ chính xác khoảng 82\% \parencite{farm2market_2025}. Dự báo nhu cầu cần dữ liệu giao dịch mà nền tảng mới chưa có nên để ở hướng phát triển. Dữ liệu mở về nông sản rất hạn chế: mục ``Nông nghiệp'' trên Cổng dữ liệu mở quốc gia chỉ có hai bộ dữ liệu \parencite{opendata_govvn}, còn số liệu của Cục Thống kê theo quý và năm \parencite{cucthongke_nongnghiep} chỉ đủ dựng đường nền theo mùa vụ. Do đó \label{subsec:nguondulieunongsan}hướng khả thi kết hợp ba lớp: số liệu vĩ mô làm đường nền, trong đó giá nhà sản xuất theo tháng của FAOSTAT \parencite{faostat_producerprices} là nguồn huấn luyện duy nhất đủ dài, còn bản tin giá tuần của Trung tâm Khuyến nông Lâm Đồng \parencite{khuyennong_lamdong_gia} mới có vài tuần liên tục nên chỉ làm bằng chứng giá trong nước; dữ liệu Farmery tích lũy (giai đoạn đầu dùng phương pháp thống kê đơn giản theo mùa vụ); hệ thống truy xuất quốc gia dùng để chuẩn hóa lược đồ định danh chứ không để lấy dữ liệu huấn luyện. Hạn chế này được nêu lại khi đánh giá ở Chương~\ref{chap:danhgia}.

\textbf{Giám sát dấu hiệu bất thường.} Các kỹ thuật học máy (thống kê, phân loại, phân cụm \parencite{chandola2009anomaly}) cần lưu lượng giao dịch lớn và tập nhãn gian lận đã xác nhận --- hai điều kiện nền tảng mới chưa có, nên không kiểm chứng được bằng số đo (mục~\ref{sec:muctieudanhgia}). Trong đồ án, Farmery dùng \emph{luật ngưỡng} do quản trị viên cấu hình (nhiều đánh giá năm sao cùng người mua trong thời gian ngắn, tỷ lệ hủy đơn vượt mức, giá lệch quá biên độ so với trung vị danh mục) làm lớp cảnh báo trước khi nhân sự xem xét; thay bằng học máy khi đủ dữ liệu được để ở hướng phát triển (Chương~\ref{chap:ketluan}).
